\documentclass[11pt]{article}

\usepackage[T1]{fontenc}
\usepackage[utf8]{inputenc}
\usepackage{lmodern}
\usepackage{geometry}
\usepackage{amsmath,amssymb,amsthm,mathtools}
\usepackage{microtype}
\usepackage{hyperref}
\usepackage{enumitem}
\usepackage{booktabs}

\geometry{margin=1in}
\hypersetup{colorlinks=true,linkcolor=blue,citecolor=blue,urlcolor=blue}

\newtheorem{definition}{Definition}[section]
\newtheorem{theorem}[definition]{Theorem}
\newtheorem{proposition}[definition]{Proposition}
\newtheorem{corollary}[definition]{Corollary}
\newtheorem{remark}[definition]{Remark}

\newcommand{\Wc}{\mathcal W_c}
\newcommand{\WT}{\mathcal W_T}
\newcommand{\Vk}{\mathcal V_T}
\newcommand{\BT}{\mathcal B_T}
\newcommand{\riskc}{\rho_c}

\title{When Statistical Certificates Cannot Justify Claims:\\
Evidence-Bounded Promotion for AI Attribution}
\author{Anonymous draft}
\date{}

\begin{document}

\maketitle

\begin{abstract}
A statistical certificate controls a particular risk functional on a declared
class of data-generating worlds. It does not automatically license a stronger
semantic claim, a claim on a broader deployment population, or a claim about a
different target functional. This paper develops a conservative
certificate-to-claim interface for reasoning about such promotion steps. The
central mathematical object is a sharp claim-risk envelope over worlds
compatible with the realized certificate bounds. We show that this envelope is
the pointwise least monotone numerical warrant on its feasible domain; practical
warrants are therefore conservative upper approximations of a canonical object,
rather than competing definitions. The operational rule is fail closed:
coverage gaps, incompatible constraints, and failed computations produce an
unresolved diagnostic rather than a safety conclusion. We use AI attribution as
a motivating application, where closed-set performance certificates can be
insufficient for claims about known membership, out-of-distribution status, or
semantic provenance. The probability algebra used by the interface is
standard; the contribution is methodological and application-driven rather
than a claim of a new probability theory or type theory.
\end{abstract}

\section{Introduction}
% Drafting target: motivate the negative result and the certificate-to-claim
% problem. Connect the RAID-style negative result to a general methodological
% question: when does a valid statistical certificate permit a downstream
% semantic claim? Do not claim a new general theory.

A recurring failure mode in statistical evaluation is to treat a valid bound
on one target as permission to assert a stronger claim. A classifier may have a
well-calibrated closed-set error bound while the downstream statement concerns
whether a sample belongs to the known family at all. A detector may be
evaluated on one population while the claim is used on another. In both cases,
the statistical certificate can remain correct while the promoted claim is
unsupported.

This paper studies the promotion step itself. The aim is not to replace
existing work on partial identification, probabilistic program logics, or
assurance reasoning. Instead, we ask for a narrow operational interface that
makes three objects explicit: the world class on which a certificate is valid,
the risk functional it actually bounds, and the downstream claim-risk
functional that must be controlled before promotion.

Our main mathematical device is a sharp claim-risk envelope over worlds
compatible with the realized certificate inequalities. This yields a canonical
least monotone warrant on its feasible domain. Operationally, the interface is
fail closed: an uncovered deployment world, an empty compatible set, or a
failed optimization is not interpreted as evidence that the claim is safe.

\section{Setup and Terminology}
% Full version will consolidate the definitions already audited in
% docs/certificate-scope-effect-system-v0.2.md and
% docs/certificate-to-claim-target-typing-candidate.md.

A statistical certificate procedure is represented by
\[
\kappa_i
=
(E_i,\mathcal W_{\kappa_i}^{A},\rho_i,B_i,\alpha_i),
\]
where \(E_i\) is the evidence channel,
\(\mathcal W_{\kappa_i}^{A}\) is the world class on which the certificate's
validity assumptions hold, \(\rho_i\) is the certified target risk,
\(B_i(Z)\) is a data-dependent upper bound, and \(\alpha_i\) is its failure
budget. Uniform frequentist validity means
\[
\inf_{D\in\mathcal W_{\kappa_i}^{A}}
\Pr_D\!\left[\rho_i(D)\le B_i(Z)\right]
\ge 1-\alpha_i.
\]

A downstream claim is represented by
\[
c=(C,\Wc,\riskc,\delta_c,\bar\alpha_c),
\]
where \(\Wc\) is the deployment world class, \(\riskc\) is the risk of
asserting the claim, \(\delta_c\) is the tolerated claim risk, and
\(\bar\alpha_c\) is the maximum allowed false-promotion budget.

The distinction between a certificate risk \(\rho_i\) and the downstream
claim risk \(\riskc\) is central: a small bound on the former does not by
itself provide a numerical bound on the latter.

\section{Sharp Envelope and Canonical Least Warrant}

\subsection{Common semantic domain}

Let \(\kappa_1,\ldots,\kappa_m\) be certificates used conjunctively to support
claim \(c\), and let \(T\) declare an additional warrant scope \(\WT\).
Define
\[
\mathcal W
=
\Wc
\cap
\WT
\cap
\bigcap_{i=1}^{m}\mathcal W_{\kappa_i}^{A}.
\]
For the conjunctive rule, operational promotion requires the fail-closed
coverage condition
\[
\boxed{
\Wc
\subseteq
\WT
\cap
\bigcap_{i=1}^{m}\mathcal W_{\kappa_i}^{A}.
}
\]
The use of an intersection is essential: every certificate invoked by the
warrant must be valid on every deployment world on which the downstream claim
is asserted.

For a candidate vector of upper bounds
\(b=(b_1,\ldots,b_m)\in[0,1]^m\), define the compatible-world set
\[
\Vk(b)
=
\left\{
D\in\mathcal W:
\rho_i(D)\le b_i
\text{ for all }i
\right\}.
\]
Define the feasible-bound domain
\[
\boxed{
\BT
=
\left\{
b\in[0,1]^m:
\Vk(b)\neq\varnothing
\right\}.
}
\]

\begin{definition}[Sharp claim-risk envelope]
For \(b\in\BT\), define
\[
\boxed{
s_T^\star(b)
=
\sup_{D\in\Vk(b)}
\riskc(D).
}
\]
For \(b\notin\BT\), the operational result is the distinguished non-numeric
state \(\bot\). No numerical comparison with \(\delta_c\) is permitted in that
branch.
\end{definition}

This partial-domain definition is deliberate. Assigning a convenient numeric
value to an empty compatible set can make a later condition such as
\(s_T^\star(b)\le\delta_c\) true for the wrong reason. The operational
interface therefore separates feasibility from risk comparison.

\subsection{Canonical form}

\begin{theorem}[Canonical Least Warrant]
Assume \(\mathcal W\neq\varnothing\).

\begin{enumerate}[label=(\roman*)]
\item If \(b\in\BT\) and \(b\le b'\) coordinatewise, then
\(b'\in\BT\) and
\[
s_T^\star(b)\le s_T^\star(b').
\]
Thus \(s_T^\star\) is coordinatewise nondecreasing on its feasible-bound
domain.

\item Let \(g:[0,1]^m\to[0,1]\) be coordinatewise nondecreasing and satisfy
\[
\forall D\in\mathcal W:
\qquad
\riskc(D)
\le
g\!\left(\rho_1(D),\ldots,\rho_m(D)\right).
\]
Then, for every \(b\in\BT\),
\[
\boxed{
s_T^\star(b)\le g(b).
}
\]
Hence \(s_T^\star\) is the pointwise least monotone numerical warrant on the
feasible-bound domain induced by the declared semantic model.
\end{enumerate}
\end{theorem}

\begin{proof}
For part (i), if \(b\le b'\), every world satisfying
\(\rho_i(D)\le b_i\) for all \(i\) also satisfies
\(\rho_i(D)\le b_i'\) for all \(i\). Therefore
\[
\Vk(b)\subseteq\Vk(b').
\]
Because \(b\in\BT\), the first set is nonempty, so the second is nonempty as
well and \(b'\in\BT\). Taking suprema of the same function \(\riskc\) over
nested nonempty sets gives
\[
s_T^\star(b)\le s_T^\star(b').
\]

For part (ii), fix \(b\in\BT\) and any \(D\in\Vk(b)\). By the defining
property of \(g\) and coordinatewise monotonicity,
\[
\riskc(D)
\le
g\!\left(\rho_1(D),\ldots,\rho_m(D)\right)
\le
g(b_1,\ldots,b_m).
\]
The right-hand side is independent of \(D\). Taking the supremum over
\(D\in\Vk(b)\) yields
\[
s_T^\star(b)\le g(b).
\]
\end{proof}

\begin{corollary}[Conservative computable warrant]
Let \(\widehat s_T\) be a computable numerical function on \(\BT\) such that
\[
\widehat s_T(b)\ge s_T^\star(b)
\qquad
\forall b\in\BT.
\]
Then every \(D\in\Vk(b)\) satisfies
\[
\riskc(D)\le\widehat s_T(b).
\]
If \(\widehat s_T\) is also coordinatewise nondecreasing on \(\BT\), it is a
tractable monotone warrant on the feasible domain.
\end{corollary}

\begin{proof}
For \(D\in\Vk(b)\),
\[
\riskc(D)
\le
\sup_{D'\in\Vk(b)}\riskc(D')
=
s_T^\star(b)
\le
\widehat s_T(b).
\]
\end{proof}

The relationship between the semantic optimum and a practical warrant is
therefore
\[
\boxed{
s_T^\star(b)
\le
\widehat s_T(b),
\qquad b\in\BT.
}
\]
The sharp envelope is canonical; a closed-form, optimization-based, or
machine-verified \(g_T\) is useful when it provides a conservative and more
tractable upper approximation.

\subsection{Operational consequence}

The theorem does not turn infeasibility into safety. Runtime promotion must
first establish \(b\in\BT\). If \(\Vk(b)=\varnothing\), the result is
\(\bot\), together with a diagnostic such as inconsistent assumptions or
inconsistent realized bounds. Only after feasibility is established is a
comparison of \(s_T^\star(b)\), or of a certified conservative upper
approximation, with \(\delta_c\) semantically meaningful.

\section{Ternary Judgment and Diagnostic Reasons}
% Planned content: PROMOTE / ABSTAIN / UNRESOLVED as verification outcomes;
% reason codes remain diagnostics, not a novelty claim and not a third
% decision-theoretic action unless a loss is explicitly defined.

\section{Fail-Closed Promotion}
% Planned content: frequentist false-promotion proposition, scope-gaps, empty
% compatible sets, and examples of fail-open bugs.

\section{Application to AI Attribution}
% Planned content: closed-set certificate -> known-membership/OOD/provenance
% non-transfer and executable finite counterworld demonstration.

\section{Relation to Prior Work}
% Planned content: partial identification; aHL; Eris; BHL; StatWhy;
% Assurance 2.0/SACM; proxy-evidence world-fiber/counterworld work.
% Explicitly separate known prior art from the remaining methodological claim.

\section{Limitations and Honest Boundaries}
This paper does not claim a new probability theorem, a new theory of partial
identification, or a new type-and-effect system. The sharp-envelope theorem is
used to put the certificate-to-claim interface in canonical mathematical form.
The intended contribution is methodological: to make scope, target mismatch,
feasibility, and fail-closed claim promotion explicit in an AI-attribution
setting.

\section{Conclusion}
% Planned closing: Probability is not permission to claim. A certificate
% licenses only the downstream claims whose scope and target are connected by
% an explicit validated warrant.

\bibliographystyle{plain}
\bibliography{references}

\end{document}
